% !TEX program = xelatex
% ============================================================
%  Arabic Non-Disclosure Agreement (NDA) Template — Nibras Center
%  اتفاقية عدم الإفصاح
%  Compile with: xelatex main.tex (twice)
% ============================================================

\documentclass[12pt,a4paper]{article}

\usepackage[a4paper,margin=2.5cm,top=2.5cm,bottom=2.5cm,headheight=15pt]{geometry}
\usepackage{array}
\usepackage{booktabs}
\usepackage{tabularx}
\usepackage{enumitem}
\usepackage{float}
\usepackage{titlesec}
\usepackage{fancyhdr}
\usepackage{setspace}

\usepackage{bidi}
\usepackage[no-math]{fontspec}
\usepackage[quiet]{polyglossia}

\setmainfont[Scale=1]{NotoKufiArabic-Regular.ttf}
\newfontfamily\arabicfont[Script=Arabic,Scale=0.95]{NotoKufiArabic-Regular.ttf}
\newfontfamily\englishfont[Scale=0.95]{TeX Gyre Termes}

\setdefaultlanguage[calendar=gregorian,hijricorrection=1]{arabic}
\setotherlanguage[variant=british]{english}

\usepackage[breaklinks,hidelinks]{hyperref}
\hypersetup{pdftitle={اتفاقية عدم الإفصاح},pdfauthor={الشركة}}

\titleformat{\section}{\large\bfseries}{\thesection}{0.7em}{}
\titlespacing*{\section}{0pt}{1.2ex plus 0.3ex}{0.8ex plus 0.2ex}

\pagestyle{fancy}
\fancyhf{}
\fancyhead[R]{\small\itshape اتفاقية عدم الإفصاح (NDA)}
\fancyhead[L]{\small اسم الشركة}
\fancyfoot[C]{\small\thepage}
\renewcommand{\headrulewidth}{0.3pt}

\newcolumntype{L}[1]{>{\raggedright\arraybackslash}p{#1}}
\newcolumntype{C}[1]{>{\centering\arraybackslash}p{#1}}

\setstretch{1.5}

\begin{document}

\begin{center}
{\LARGE\bfseries اتفاقية عدم الإفصاح وحماية المعلومات}\\[0.5em]
{\large (\LR{Non-Disclosure Agreement — NDA})}\\[1em]
\end{center}

\vspace{1em}

% ============================================================
% PARTIES
% ============================================================
\noindent\textbf{الأطراف:}

\noindent وقّع الطرفان أدناه على هذه الاتفاقية:

\begin{enumerate}
    \item \textbf{الطرف الأول (المُفصِح):} \rule{8cm}{0.4pt}، الشركة المسجّلة في \rule{4cm}{0.4pt}، رقم السجل التجاري: \rule{4cm}{0.4pt}، ويمثلها في التوقيع \rule{6cm}{0.4pt}.

    \item \textbf{الطرف الثاني (المُستلِم):} \rule{8cm}{0.4pt}، رقم الهوية/الجواز: \rule{4cm}{0.4pt}، العنوان: \rule{6cm}{0.4pt}.
\end{enumerate}

\vspace{1em}

% ============================================================
% 1. PURPOSE
% ============================================================
\section{الغرض من الاتفاقية}
\label{sec:purpose}
يملك الطرف الأول معلومات سرية تتعلق بأعماله ومنتجاته وتقنياته. يرغب الطرفان في تبادل معلومات لأغراض \rule{8cm}{0.4pt} (المشروع/التعاون)، ويلتزم الطرف الثاني بالحفاظ على سرية هذه المعلومات وفقاً لشروط هذه الاتفاقية.

% ============================================================
% 2. DEFINITION
% ============================================================
\section{تعريف المعلومات السرية}
\label{sec:definition}
تشمل «المعلومات السرية» جميع المعلومات التي يُفصح عنها الطرف الأول للطرف الثاني، سواء كانت مكتوبة أو شفهية أو إلكترونية، وتشمل على سبيل المثال لا الحصر:
\begin{enumerate}
    \item الخطط التجارية واستراتيجيات العمل.
    \item الأسرار التجارية والمعرفة الفنية (\LR{Know-how}).
    \item الكود المصدري والتصاميم التقنية.
    \item قوائم العملاء والموردين.
    \item البيانات المالية والتوقعات.
    \item الملكية الفكرية وبراءات الاختراع.
    \item أي معلومات أخرى يُعلم الطرف الأول الطرف الثاني بأنها سرية.
\end{enumerate}

% ============================================================
% 3. OBLIGATIONS
% ============================================================
\section{التزامات الطرف المستلم}
\label{sec:obligations}
يلتزم الطرف الثاني بما يلي:
\begin{enumerate}
    \item \textbf{السرية:} الحفاظ على سرية المعلومات وعدم إفشائها لأي طرف ثالث.
    \item \textbf{الاستخدام المحدود:} استخدام المعلومات لأغراض الغرض المذكور في البند 1 فقط.
    \item \textbf{الحماية:} اتخاذ نفس درجة الحماية التي يتخذها لحماية معلوماته السرية الخاصة (ولكن لا تقل عن درجة معقولة).
    \item \textbf{الوصول المحدود:} عدم السماح بالوصول للمعلومات إلا للأشخاص الذين يحتاجونها لأغراض الغرض المذكور، والذين وقّعوا اتفاقيات سرية مماثلة.
    \item \textbf{عدم النسخ:} عدم نسخ أو توثيق المعلومات إلا بقدر الضرورة.
\end{enumerate}

% ============================================================
% 4. EXCLUSIONS
% ============================================================
\section{الاستثناءات}
\label{sec:exclusions}
لا تشمل هذه الاتفاقية المعلومات التي:
\begin{enumerate}
    \item كانت معروفة للطرف الثاني قبل تاريخ الإفصاح.
    \item أصبحت عامة دون إخلال من الطرف الثاني.
    \item حصل عليها الطرف الثاني بشكل مشروع من طرف ثالث دون التزام بالسرية.
    \item طوّرها الطرف الثاني بشكل مستقل دون استخدام المعلومات السرية.
    \item طُلبت للإفصاح بموجب أمر قضائي أو حكم قانوني، بشرط إشعار الطرف الأول مسبقاً.
\end{enumerate}

% ============================================================
% 5. DURATION
% ============================================================
\section{مدة السرية}
\label{sec:duration}
\begin{enumerate}
    \item \textbf{مدة الاتفاقية:} تبدأ من تاريخ التوقيع وتستمر لمدة \rule{2cm}{0.4pt} سنوات.
    \item \textbf{مدة الالتزام بالسرية:} يستمر الالتزام بالسرية لمدة \rule{2cm}{0.4pt} سنوات بعد انتهاء الاتفاقية.
    \item \textbf{الأسرار التجارية:} تبقى الأسرار التجارية سرية إلى الأبد ما لم تصبح عامة.
\end{enumerate}

% ============================================================
% 6. RETURN
% ============================================================
\section{إرجاع أو إتلاف المعلومات}
\label{sec:return}
عند انتهاء الغرض من الاتفاقية أو عند طلب الطرف الأول، يلتزم الطرف الثاني بـ:
\begin{enumerate}
    \item إرجاع جميع الوثائق والمواد التي تحتوي على معلومات سرية.
    \item أو إتلافها بشكل آمن وتقديم شهادة مكتوبة بالإتلاف.
    \item عدم الاحتفاظ بأي نسخ أو نُسخ إلكترونية.
\end{enumerate}

% ============================================================
% 7. OWNERSHIP
% ============================================================
\section{ملكية المعلومات}
\label{sec:ownership}
\begin{enumerate}
    \item تبقى جميع المعلومات السرية ملكاً حصرياً للطرف الأول.
    \item لا تمنح هذه الاتفاقية الطرف الثاني أي حق أو ترخيص في المعلومات.
    \item لا تُعتبر هذه الاتفاقية تنازلاً عن أي حق من حقوق الملكية الفكرية.
\end{enumerate}

% ============================================================
% 8. REMEDIES
% ============================================================
\section{العقوبات والتعويض}
\label{sec:remedies}
\begin{enumerate}
    \item يقرّ الطرف الثاني أن الإخلال بالسرية قد يسبب أضراراً غير قابلة للتعويض المالي.
    \item يحق للطرف الأول الحصول على أحكام قضائية عاجلة (\LR{Injunction}) لمنع الإخلال.
    \item يحق للطرف الأول المطالبة بالتعويض عن جميع الأضرار الناتجة عن الإخلال.
\end{enumerate}

% ============================================================
% 9. GOVERNING LAW
% ============================================================
\section{القانون الحاكم وحل النزاعات}
\label{sec:law}
\begin{enumerate}
    \item تُفسَّر هذه الاتفاقية وفقاً لقوانين \rule{6cm}{0.4pt}.
    \item تُحلّ النزاعات عن طريق التحكيم في \rule{4cm}{0.4pt}.
\end{enumerate}

% ============================================================
% SIGNATURES
% ============================================================
\vspace{2em}
\section*{التوقيعات}

\begin{table}[H]
\centering
\renewcommand{\arraystretch}{2.5}
\begin{tabularx}{\textwidth}{L{6cm} X X}
\toprule
& \textbf{الطرف الأول} & \textbf{الطرف الثاني} \\
\midrule
\textbf{الاسم} & --- & --- \\
\textbf{الصفة} & --- & --- \\
\textbf{التوقيع} & \rule{4cm}{0.4pt} & \rule{4cm}{0.4pt} \\
\textbf{التاريخ} & \rule{4cm}{0.4pt} & \rule{4cm}{0.4pt} \\
\bottomrule
\end{tabularx}
\end{table}

\end{document}
